\documentclass[11pt,a4paper]{article}

\usepackage[utf8]{inputenc}
\usepackage[T1]{fontenc}
\usepackage{amsmath,amssymb,amsthm}
\usepackage{graphicx}
\usepackage{booktabs}
\usepackage{hyperref}
\usepackage[margin=2.5cm]{geometry}

\title{A Casimir-Locked Bosonic-Seesaw Ansatz for the Electroweak Boson Masses}
\author{[Author]\\
\small [Affiliation]}
\date{16 May 2026 -- v4 (PRD submission candidate)}

\newcommand{\md}{\mathbf{2}}
\newcommand{\mt}{\mathbf{3}}
\newcommand{\Cas}{C_2}
\newcommand{\bsm}{BSM}

\begin{document}

\maketitle

\begin{abstract}
We propose a Beyond-Standard-Model ansatz for the electroweak boson mass
spectrum in which the quadratic mass terms of the scalar seed sector are
\emph{locked} by the SU(2) Casimir invariants of two internal channels: a
doublet ($\md$) and a triplet ($\mt$, adjoint). For each channel we
postulate a $2\times 2$ mass-squared seed matrix in the
elementary/composite basis with a protected zero on the diagonal, an
off-diagonal entry $\sqrt{\Cas(R)}\,m_0$, and a tachyonic diagonal entry
$-\Cas(R)\, m_0^2$. Diagonalisation gives four closed-form eigenvalues.
We then \emph{postulate} the identifications
$M_W^2 = m_0^2\,\lambda_{\md,+}$, $M_Z^2 = m_0^2\,\lambda_{\mt,+}$,
$m_{h,{\rm bare}}^2 = m_0^2\,|\lambda_{\md,-}|$, and
$(v/\sqrt{2})^2 = m_0^2\,|\lambda_{\mt,-}| - \Delta$. A single
bidirectional spurion of magnitude $\Delta = (26.6~\text{GeV})^2$ acts
in opposite directions on the doublet and triplet tachyonic
eigenvalues. With $M_W$ as the anchor (fixing $m_0 = 106.57$~GeV) and
either $m_h$ or $G_F$ anchoring $\Delta$, the model parameter-freely
predicts $M_Z = 91.18$~GeV (PDG: 91.19, dev $0.007\%$) and
$m_{h,{\rm bare}} = 122.39$~GeV. The remaining EW observable in
$\{m_h, v/\sqrt{2}\}$ is matched to $0.04\%$ by the cross-irrep
negative-branch trace identity $m_h^2 + (v/\sqrt{2})^2 =
m_0^2(|\lambda_{\md,-}| + |\lambda_{\mt,-}|)$ at the
$(5~\text{GeV})^2$ level. The doublet half of the spurion $\Delta$
is reproduced to $1\,\%$ by the one-loop top Coleman--Weinberg
shift at cutoff $m_0$~\cite{ColemanWeinberg1973, Haba2016BosonicSeesaw};
the triplet half receives only a $\sim\!15\times$-too-small contribution
from gauge loops alone, leaving the bidirectional picture's triplet
shift as an empirical feature without a natural one-loop mechanism. A representation-selection
analysis (Sec.~\ref{sec:repsel}) shows the doublet--triplet pair is the
unique numerical winner among the ten SU(2) irrep pairs with
$\Cas < 12$, by a factor of $\sim 300$ in maximum fractional error. As
a structural side-prediction, extending the seed to the next internal
irrep ($R = 4$, $j = 3/2$) generates a scalar at $96.5$~GeV; this lies
in the mass range of the CMS+LEP $\gamma\gamma$ hint near $95.4$~GeV
but, at $1.2\,\%$ deviation, is $\sim 30\times$ looser than the
$M_W, M_Z, m_h, v$ predictions of this paper, so we do not claim a
quantitative match. We do \emph{not} yet derive the
eigenvalue-to-physical-mass identifications from a vacuum-expectation
mechanism --- the gauge mass matrix and the scalar potential
minimisation remain open work, listed explicitly in
Section~\ref{sec:status}. The framework is positioned within the
bosonic-seesaw, partial-compositeness, and top-seesaw literatures.
A complete computational pipeline (SymPy + a hand-written UFO +
MadGraph) reproduces every numerical statement in this paper.
\end{abstract}

\section{Introduction}

The Standard Model gauge group $SU(2)_L \times U(1)_Y$ accounts for the
observed electroweak boson masses through the Higgs mechanism, but treats
the scalar potential parameters $\mu^2$, $\lambda$ as free. Decades of
beyond-Standard-Model (BSM) work have sought to derive these parameters
from a deeper structure: composite Higgs, technicolor, supersymmetric
fixed points, top-seesaw, asymptotic safety, gauge-Higgs unification, and
many more. Sch\"ucker's pre-discovery
compilation~\cite{Schucker2007} catalogues roughly one hundred
predictions for $m_h$.

In this work we follow a structurally different route. We \emph{do not}
attempt to derive $m_h$, $M_W$, or $M_Z$ from a UV mechanism in the
conventional sense (Higgs-mechanism minimisation of an explicit potential
followed by gauge-boson mass extraction). Instead, we ask whether a
purely numerical ansatz --- a $2\times 2$ scalar mass-squared matrix
whose entries are locked by SU(2) Casimirs --- can reproduce the
electroweak boson spectrum once specific eigenvalues are
\emph{identified} with the physical masses. We find that with $M_Z$ and
$G_F$ as inputs, this identification reproduces $M_W$ and $m_h$ at the
$10^{-3}$ level.

This is a numerologically suggestive ansatz, not a closed derivation. The
gauge-boson mass matrix is not derived from VEVs in the present draft;
the identification $M_W^2 = m_0^2 \lambda_{\md,+}$,
$M_Z^2 = m_0^2 \lambda_{\mt,+}$ is a postulate. Likewise the Higgs-mass
formula is asserted on top of the postulated eigenvalue assignment, with
the quartic couplings $\lambda_d, \lambda_t, \lambda_{dt}$ left free.
Section~\ref{sec:status} states these open derivations explicitly.

The mechanism is a representation-theoretic version of the
\emph{bosonic seesaw} of Calmet--Oliver~\cite{Calmet2006} and
Kim~\cite{Kim2005}: a $2\times 2$ scalar mass-squared matrix whose
diagonalisation produces one tachyonic eigenvalue, structurally
suggestive of EWSB. What is new is the \emph{Casimir locking} of the
matrix entries.

\section{Literature landscape}

Table~\ref{tab:landscape} compares the present model to its closest
relatives in the literature.

\begin{table}[ht]
\centering
\small
\begin{tabular}{p{3.6cm} p{2.6cm} p{2.6cm} p{3.0cm} p{3.0cm}}
\toprule
\textbf{Framework} & \textbf{Mechanism} & \textbf{Gauge group} &
\textbf{Scalar content} & \textbf{Distinctive feature} \\
\midrule
Standard Model & Higgs mechanism & $SU(2)_L \times U(1)_Y$ &
1 doublet ($Y=1/2$) & $\mu^2$, $\lambda$ free \\
Calmet--Oliver~\cite{Calmet2006} & Bosonic seesaw & $SU(2)_L \times U(1)_Y$ &
2 Higgs doublets & EWSB from negative det \\
Kim~\cite{Kim2005} & SUSY + bosonic seesaw & MSSM + $\mt$ &
2 doublets + triplet & EWSB tied to SUSY-breaking \\
Caracciolo--Parolini--Serone~\cite{Caracciolo2012} & Composite Higgs &
$SU(2)_L \times U(1)_Y$ + $SO(N)_{HC}$ & pNGB doublet from coset &
Seiberg-dual UV completion \\
Chivukula \emph{et al.}~\cite{Chivukula1998} & Top-seesaw &
$SU(2)_L \times U(1)_Y$ & composite Higgs from $t_L$ &
fermionic (not bosonic) seesaw \\
Bando--Kugo--Yamawaki~\cite{Bando1985} & Hidden local symmetry &
emergent local symmetry & vector mesons as gauge bosons & KSRF relations \\
\midrule
\textbf{This work} & Casimir-locked bosonic seesaw &
$SU(2)_L \times U(1)_Y$ & $(\Phi^d_e, \Phi^d_c, \Phi^t_e, \Phi^t_c)$ &
$M^2_R / m_0^2$ \emph{fixed by} $\Cas(R)$ \\
\bottomrule
\end{tabular}
\caption{Position of the present model within the BSM landscape.}
\label{tab:landscape}
\end{table}

The closest structural precedent is Kim's SUSY bosonic-seesaw with weak
triplet~\cite{Kim2005}: he uses the same matrix architecture (bosonic
seesaw), the same field content (doublet + triplet), and the same EWSB
mechanism (negative determinant). What we add is the
\emph{representation-theoretic constraint} on the matrix entries.

\section{The Casimir-locked seed}

Introduce a single mass scale $m_0$ and assume that the EWSB sector
contains two protected internal channels carrying the SU(2) irreps $\md$
and $\mt$. For each irrep $R$, introduce two scalars $\Phi^R_e, \Phi^R_c$
in the \emph{elementary/composite basis} of partial
compositeness~\cite{Caracciolo2012}, with the quadratic seed
\begin{equation}
  \mathcal L^R_{\rm seed} \;=\;
  - \begin{pmatrix} \Phi^{R\dagger}_e & \Phi^{R\dagger}_c \end{pmatrix}
  \mathcal M_R^2
  \begin{pmatrix} \Phi^{R}_e \\ \Phi^{R}_c \end{pmatrix},
  \qquad
  \frac{\mathcal M_R^2}{m_0^2}
  \;=\;
  \begin{pmatrix}
    0 & a\,\sqrt{\Cas(R)} \\
    a\,\sqrt{\Cas(R)} & -\,b\,\Cas(R)
  \end{pmatrix}.
\label{eq:seed}
\end{equation}
The Casimir is normalised by $T^a T^a = \Cas(R)\,\mathbf{1}_R$ with
generators obeying $[T^a, T^b] = i\epsilon^{abc} T^c$; the doublet has
$T^a = \sigma^a/2$ giving $\Cas(\md) = 3/4$, and the adjoint triplet has
$\Cas(\mt) = 2$. Three structural conditions specify
Eq.~\eqref{eq:seed} up to the dimensionless coefficients $a, b$:
(i) one diagonal entry is set to zero, conventionally the elementary
$(1,1)$ slot, by a gauge or discrete symmetry that protects it from a
bilinear mass; (ii) the trace polynomial in $\Cas(R)$ has no constant
term and minimal Casimir power, fixing the $(2,2)$ entry as $-b\,\Cas(R)$;
(iii) the off-diagonal mixing scales with the SU(2) generator norm
$\sqrt{\Cas(R)} = t_R$, the lowest odd power of the Casimir consistent
with the group-theoretic content.

\paragraph{Status of $a = b = 1$.} The coefficient $a$ can be absorbed
into a rescaling of $m_0$ by canonically normalising the kinetic term of
the protected $(1,1)$ field; setting $a = 1$ is therefore a convention,
not a prediction. The coefficient $b$ has no analogous reparameterisation
and is, at present, an independent postulate: $b = 1$ is the simplest
non-trivial integer Wilson coefficient in the polynomial expansion of
$m_{22}$ in $\Cas$, and we have not found a symmetry argument
(custodial-SU(2) Ward identity, trace condition, or kinetic
normalisation) that uniquely fixes it. The ``Casimir-locked'' label
refers to the \emph{shape} of Eq.~\eqref{eq:seed} -- both entries scale
with the same $\Cas(R)$ with no $R$-dependent coefficients --
not to a derivation of $b$. We work with $a = b = 1$ throughout.

For SU(2), $\Cas(\md) = 3/4$ and $\Cas(\mt) = 2$. The two seed blocks are
\begin{equation}
  \frac{\mathcal M_{\md}^2}{m_0^2} = \begin{pmatrix} 0 & \sqrt{3}/2 \\
  \sqrt{3}/2 & -3/4 \end{pmatrix},
  \qquad
  \frac{\mathcal M_{\mt}^2}{m_0^2} = \begin{pmatrix} 0 & \sqrt{2} \\
  \sqrt{2} & -2 \end{pmatrix}.
\end{equation}
Their characteristic equations $\lambda^2 + \tfrac{3}{4}\lambda - \tfrac{3}{4} = 0$
and $\lambda^2 + 2\lambda - 2 = 0$ have closed-form roots
\begin{equation}
  \lambda_{\md,\pm} = \frac{-3 \pm \sqrt{57}}{8},
  \qquad
  \lambda_{\mt,\pm} = -1 \pm \sqrt{3},
\label{eq:eigs}
\end{equation}
giving four protected mass-squared scales in units of $m_0^2$:
\[
  +\,\tfrac{\sqrt{57}-3}{8},\quad
  -\,\tfrac{\sqrt{57}+3}{8},\quad
  +\,(\sqrt{3}-1),\quad
  -\,(\sqrt{3}+1).
\]

\paragraph{Robustness in $(a, b)$.} A grid scan of the seed coefficients
(see \texttt{algebra/robustness.py}) decomposes naturally into two
directions. Along $a = b$ the model is gauge-trivial: the rescaling is
absorbed into $m_0$ and predictions are unchanged. Orthogonal to this
direction (the $a/b$ ratio), the model is fine-tuned at the $\sim 10^{-3}$
level: $\partial(\delta m_h / m_h)/\partial(a/b)\big|_{(1,1)}
\simeq +0.67$, so a $1\%$ misalignment of $a$ vs $b$ shifts $m_h$ by
$\sim 0.67\%$. The $0.04\%$ accuracy on $m_h$ requires $a/b = 1$ to
$\sim\!6 \times 10^{-4}$. The Casimir lock is therefore the operational
postulate that the off-diagonal mixing $\sqrt{\Cas(R)}$ and the diagonal
$-\Cas(R)$ carry the \emph{same} prefactor; no symmetry argument
deriving this is given here. A natural origin is sketched in
Section~\ref{sec:PCsketch}.

\section{Representation selection: why the doublet--triplet pair?}
\label{sec:repsel}

The seed Eq.~\eqref{eq:seed} is parameterised by the SU(2) irrep $R$,
but the model so far does not say which irreps to use. Two natural
selection rules are available, one parameter-free and one numerical;
they coincide.

\paragraph{Rule A (group-theoretic):} use the two lowest non-trivial
irreducible representations of SU(2), i.e.\ the dimension-2 doublet and
the dimension-3 triplet. This rule is parameter-free and uses no
phenomenological input.

\paragraph{Rule B (numerical):} enumerate all unordered pairs
$\{R_1, R_2\}$ of SU(2) irreps with dimension $n \in \{2, 3, 4, 5, 6\}$
(equivalently $\Cas < 12$), anchor $m_0$ by setting the lighter
positive eigenvalue equal to $M_W^2$, then test the remaining three
eigenvalues against $\{M_Z, m_h, v/\sqrt{2}\}$. Table~\ref{tab:repsel}
records the maximum fractional deviation across $\{M_Z, m_h\}$.

\begin{table}[ht]
\centering
\small
\begin{tabular}{c c c c}
\toprule
pair & $(\Cas^{(1)}, \Cas^{(2)})$ & best $m_0$ (GeV) & max frac.\ error \\
\midrule
$\{2,3\}$ & $(0.75, 2.00)$ & 106.57 & \textbf{0.02 \%} \\
$\{2,4\}$ & $(0.75, 3.75)$ & 106.57 & 5.86 \% \\
$\{2,5\}$ & $(0.75, 6.00)$ & 106.57 & 9.19 \% \\
$\{2,6\}$ & $(0.75, 8.75)$ & 106.57 & 11.25 \% \\
$\{3,4\}$ & $(2.00, 3.75)$ & 88.73  & 19.82 \% \\
$\{3,5\}$ & $(2.00, 6.00)$ & 86.02  & 19.29 \% \\
$\{3,6\}$ & $(2.00, 8.75)$ & 84.43  & 20.78 \% \\
$\{4,5\}$ & $(3.75, 6.00)$ & 86.02  & 50.24 \% \\
$\{4,6\}$ & $(3.75, 8.75)$ & 84.43  & 47.46 \% \\
$\{5,6\}$ & $(6.00, 8.75)$ & 84.43  & 80.83 \% \\
\bottomrule
\end{tabular}
\caption{Ten-pair representation enumeration. For each unordered pair of
SU(2) irreps with $\Cas < 12$, $m_0$ is fitted to the lightest-positive
eigenvalue (anchoring $M_W$); the right column is the maximum
fractional deviation across $\{M_Z, m_h\}$. The doublet--triplet pair
$\{2,3\}$ is the unique numerical winner by a factor of $\sim 300$.}
\label{tab:repsel}
\end{table}

The dominance of $\{2,3\}$ has two independent algebraic sources.
First, only the doublet ($n = 2$) delivers
$\sqrt{|\lambda_-/\lambda_+|} \approx 1.523$, the ratio that maps
$M_W = 80.369$~GeV to $m_{h,{\rm bare}} \approx 122.4$~GeV; higher $n$
overshoots monotonically. Second, with $R_1 = 2$ fixed, the requirement
$\sqrt{\lambda_+(R_2) / \lambda_+(2)} = M_Z / M_W = 1.1346$ is satisfied
only by $\lambda_+(R_2) = \sqrt{3} - 1$, which singles out the triplet.
Rules A and B therefore coincide: the doublet--triplet pair is selected
both as the two lowest irreps and as the unique numerical winner.

\section{Lagrangian and gauge group}

At the EW scale we assume the Standard Model gauge group $SU(2)_L \times
U(1)_Y$ unchanged: hypercharge then acts as a low-energy spurion that
splits the doublet- and triplet-channel mass eigenstates by $Y$,
consistent with the observed $W$/$Z$/$\gamma$ spectrum. The Lagrangian is
\begin{align}
  \mathcal L &= -\frac14 W^a_{\mu\nu}W^{a\,\mu\nu} - \frac14 B_{\mu\nu}B^{\mu\nu}
  + \mathcal L_{\rm kin}^{\rm scalar} - V,\\
  V &= V_{\rm seed}^{\md} + V_{\rm seed}^{\mt}
       + \lambda_d (\Phi^{d\dagger}_c \Phi^d_c)^2
       + \lambda_t (\Phi^t_c \cdot \Phi^t_c)^2
       + \lambda_{dt} (\Phi^{d\dagger}_c \Phi^d_c)(\Phi^t_c \cdot \Phi^t_c),
\end{align}
with $V_{\rm seed}^R$ from Eq.~\eqref{eq:seed} and the doublet covariant
derivative
\(D_\mu \Phi^{d}_\alpha = (\partial_\mu - i g W^a_\mu T^a - i\tfrac{g'}{2} B_\mu)\Phi^{d}_\alpha\),
$\alpha \in \{e, c\}$. The triplets are taken to be real, $Y = 0$ SU(2)
triplets (this is \emph{not} the Georgi--Machacek scalar sector, which
uses a larger custodial-preserving multiplet content; we use only a
single real $Y=0$ triplet, of the family studied
in~\cite{Cirelli2006, FischerVanDerBij2014, Khan2018Triplet, Katayose2021TripletDM}). The portal coupling $\lambda_{dt}$ is
required to be small so that any induced triplet vacuum expectation
value remains below the $\rho$-parameter bound; this is an external
phenomenological constraint, not a derivation. The quartic couplings
$\lambda_d, \lambda_t, \lambda_{dt}$ are otherwise free in this draft;
see Section~\ref{sec:status}.

\section{Postulated identifications and resulting numerical agreements}
\label{sec:postulates}

Before stating the identifications, we recall a remark that Higgs himself
made in the augmentation of his second 1964 paper, the version eventually
accepted by Phys.~Rev.~Lett.~\cite{Higgs1964} and reproduced verbatim in
his 2013 Nobel lecture~\cite{HiggsNobel2014}:

\begin{quote}
``\emph{It is worth noting that an essential feature of this type of theory is the prediction of incomplete multiplets of scalar and vector bosons.}''
\end{quote}

This was Higgs's own flag, in 1964, that the symmetry-breaking mechanism
generically delivers \emph{incomplete} scalar+vector multiplets. The
historical Standard Model construction by Weinberg~\cite{Weinberg1967}
selected exactly one SU(2)$_L$ doublet (Y=1/2) -- four real degrees of
freedom, three eaten by $W^\pm, Z$, one physical Higgs -- because that
is the minimal field content satisfying $\rho = 1$ at tree level via
custodial symmetry. The triplet was excluded, on phenomenological
grounds, as a VEV-acquiring field. The Casimir-locked picture of
Section~\ref{sec:repsel} is, in this respect, the missing piece of
Higgs's ``incomplete multiplet'' remark: SU(2)$_L$ admits the entire
representation tower $\{\md, \mt, \mathbf{4}, \mathbf{5}, \dots\}$, and
the standard construction \emph{truncates} to $\md$ for empirical
($\rho$-parameter) reasons. Our seed structure makes the truncation
explicit by retaining $\mt$ as a non-VEV mass-locked sector whose
positive eigenvalue identifies with $M_Z$, with $\mathbf{4}$ generating
a parameter-free side-prediction at $96.54$~GeV
(Section~\ref{sec:j32}). The identifications
below should be read in that spirit.

\paragraph{Gauge sector.} We \emph{postulate} the identification
\begin{equation}
  M_W^2 \;=\; m_0^2\, \lambda_{\md,+},
  \qquad
  M_Z^2 \;=\; m_0^2\, \lambda_{\mt,+}.
\label{eq:gauge_id}
\end{equation}
Eq.~\eqref{eq:gauge_id} is \emph{not} derived from the Lagrangian written
above; deriving it from a Higgs-mechanism mass matrix is open work, see
Section~\ref{sec:status}. Taking it as an ansatz, the ratio gives
\begin{equation}
  \sin^2\theta_W \;=\; 1 - \frac{M_W^2}{M_Z^2}
  \;=\; 1 - \frac{\lambda_{\md,+}}{\lambda_{\mt,+}}
  \;=\; 1 - \frac{(\sqrt{57}-3)(\sqrt{3}+1)}{16}
  \;\approx\; 0.22310.
\label{eq:sin2thW}
\end{equation}
This algebraic ratio coincides with the on-shell PDG value
$\sin^2\theta_W = 1 - M_W^2/M_Z^2 \simeq 0.22290(28)$ within roughly $1$
standard deviation; we make no claim about $10^{-10}$-precision
experimental matching. With $M_Z = 91.1880$~GeV
(PDG~\cite{PDG2024}) as input, the ansatz fixes
$m_0 = M_Z/\sqrt{\sqrt{3}-1} = 106.578$~GeV and outputs
$M_W = 80.375$~GeV, to be compared with the PDG value $80.3692(11)$ GeV
(deviation $0.007\%$). Since $M_Z$ was used to set $m_0$, the
non-trivial output here is $M_W$ alone --- the ratio $M_W/M_Z$ is the
single new numerical statement.

\paragraph{Bare Higgs sector.} We further \emph{postulate}
\begin{equation}
  m_{h,{\rm bare}}^2 \;=\; |\lambda_{\md,-}|\, m_0^2,
\label{eq:higgs_bare}
\end{equation}
identifying the doublet's tachyonic eigenvalue with the bare Higgs
mass-squared. With $|\lambda_{\md,-}| = (\sqrt{57}+3)/8$ and
$m_0 = 106.578$~GeV, this gives
\begin{equation}
  m_{h,{\rm bare}} \;=\; m_0 \sqrt{\tfrac{\sqrt{57}+3}{8}} \;=\; 122.38~\text{GeV},
\end{equation}
short of the PDG value $125.20(11)$~GeV by $\Delta m_h = 2.82$~GeV. We
emphasise that Eq.~\eqref{eq:higgs_bare} is again not derived from the
scalar-potential minimisation of Section~5; the quartic couplings
$\lambda_d, \lambda_t, \lambda_{dt}$ enter the physical Higgs mass via
the standard 2HDM-plus-triplet algebra, and the relation between those
couplings and Eq.~\eqref{eq:higgs_bare} is not established here. The
triplet's tachyonic eigenvalue $|\lambda_{\mt,-}|\,m_0^2$ corresponds to
a separate heavy scalar at $M(\Phi^t_-) = 176.16$~GeV in this
identification, not to a piece of the Higgs.

\paragraph{Distinctive triplet prediction.} By construction
Eq.~\eqref{eq:gauge_id} forces a scalar seed eigenvalue at
$M(\Phi^t_+) = m_0 \sqrt{\sqrt{3}-1} = M_Z = 91.188$~GeV. If the seed
eigenstate is realised as a physical neutral scalar, the model predicts a
state mass-degenerate with the $Z$ --- a sharp falsifiable signature.
Existing LEP $Z'/H'$ searches and LHC scalar searches constrain its
couplings to fermions/bosons (open question Q4 in
Section~\ref{sec:status}).

\section{The spurion $\Delta$: bidirectional, cross-irrep, eigenstate-aligned}
\label{sec:spurion}

The bare Higgs from Eq.~\eqref{eq:higgs_bare} sits at $122.38$~GeV,
short of $m_{h,{\rm phys}} = 125.20$~GeV by $\Delta_d \equiv
m_{h,{\rm phys}}^2 - m_{h,{\rm bare}}^2 = (26.38~\text{GeV})^2$. The
triplet's tachyonic eigenvalue gives $\sqrt{|\lambda_{\mt,-}|}\,m_0 =
176.16$~GeV, larger than the Fermi-scale $v/\sqrt{2} = 174.10$~GeV
(from $G_F$) by $\Delta_t \equiv |\lambda_{\mt,-}|\,m_0^2 - (v/\sqrt{2})^2
= (26.85~\text{GeV})^2$. The two shifts are equal in magnitude to within
$3.6\,\%$,
\[
  \frac{\Delta_t}{\Delta_d} \;=\; 1.036,
\]
suggesting a single underlying spurion of magnitude $\Delta \approx
(26.6~\text{GeV})^2$ acting bidirectionally across the negative-branch
of the doublet and triplet sectors.

\paragraph{Bidirectional spurion ansatz.} We \emph{postulate} a single
mass-squared shift that lifts the doublet's tachyonic eigenvalue by
$+\Delta$ and lowers the triplet's by $-\Delta$:
\begin{equation}
  |\lambda_{\md,-}|\,m_0^2 \;+\; \Delta \;=\; m_{h,{\rm phys}}^2,\qquad
  |\lambda_{\mt,-}|\,m_0^2 \;-\; \Delta \;=\; (v/\sqrt{2})^2.
\label{eq:bidir}
\end{equation}
Adding the two equations gives the cross-irrep negative-branch trace
identity
\begin{equation}
  m_h^2 \;+\; (v/\sqrt{2})^2 \;=\; m_0^2\,\bigl(|\lambda_{\md,-}| + |\lambda_{\mt,-}|\bigr),
\label{eq:trace_id}
\end{equation}
which is satisfied by the model and PDG inputs to within $(5~\text{GeV})^2$
(i.e.\ $\sim 25~\text{GeV}^2$, $\sim 3.4\,\%$ of $\Delta$). With
$\Delta = \tfrac12(\Delta_d + \Delta_t) = (26.62~\text{GeV})^2$, the
\emph{two-parameter} model $(m_0, \Delta)$ outputs four electroweak
quantities
\[
  M_W = 80.37~\text{GeV},\quad M_Z = 91.19~\text{GeV},\quad
  m_h = 125.25~\text{GeV},\quad v/\sqrt{2} = 174.14~\text{GeV},
\]
all within $0.04\,\%$ of their PDG / $G_F$ values. The Fermi scale $v$,
previously taken as an external input, is now a model output. The
non-trivial ratio between the triplet and doublet shifts, $\Delta_t /
\Delta_d \neq 1$ at $3.6\,\%$, is a small structural residue not yet
explained.

\paragraph{Caution: per-irrep trace-preserving is excluded.} The
``bidirectional, cross-irrep'' label is precise and important. A
literal \emph{per-irrep} trace-preserving shift inside the doublet
sector, $\delta\lambda_+ + \delta\lambda_- = 0$ with
$\delta\lambda_- = -(26.42~\text{GeV})^2$, would also shift
$\delta\lambda_+ = +(26.42~\text{GeV})^2$ and move $M_W$ from
$80.37$~GeV to $84.6$~GeV --- a $4.23$~GeV mismatch,
$\sim\!\!300\sigma$ outside any plausible PDG / CMS range
\cite{CMS2026WMass}. The shift in Eq.~\eqref{eq:bidir} preserves the
trace of the negative-eigenvalue \emph{pair across the two irreps}, not
the trace of either single $\mathcal{M}_R^2$. In matrix language, it is
a rank-1 eigenstate-aligned perturbation on each irrep separately,
\begin{equation}
  \delta \mathcal M_{\md}^2 = -\Delta\, P^{(\md)}_-,\qquad
  \delta \mathcal M_{\mt}^2 = +\Delta\, P^{(\mt)}_-,\qquad
  P^{(R)}_- = |v_-^{(R)}\rangle\langle v_-^{(R)}|,
\label{eq:rank1}
\end{equation}
with opposite signs on the two channels. Each $P^{(R)}_-$ projects onto
the tachyonic eigenvector and leaves $\lambda_+(R) = M_{W,Z}^2$
untouched. The earlier ``antisymmetric splitting'' /
``trace-preserving spurion'' wording from the chat draft is replaced by
this precise statement.

\paragraph{Mechanism for the doublet shift: top-loop Coleman--Weinberg.}
The quadratically-divergent top-loop contribution to the Higgs
mass-squared, regulated by a hard cutoff at the seed scale
$\Lambda = m_0$~\cite{ColemanWeinberg1973}, evaluates to
\begin{equation}
  \Delta_{{\rm CW},\,{\rm quad}} \;=\;
    -\,\frac{3 y_t^2}{8\pi^2}\,\big(m_0^2 - m_t^2\big)
    \;=\; (26.28~\text{GeV})^2,
\label{eq:cw_top}
\end{equation}
using $y_t = \sqrt{2}\,m_t / v = 0.99196$, $m_0 = 106.57$~GeV, and
$m_t = 172.7$~GeV. This is $0.989$ of $\Delta_d = (26.38~\text{GeV})^2$
--- agreement at the $1\,\%$ level with no additional tunable scale.
The cutoff at which the top loop produces $\Delta_d$ of the right size
is \emph{exactly} the seed scale $m_0$ that Eq.~\eqref{eq:gauge_id}
fixed from the gauge anchor. Whether this is coincidence or a hint that
$m_0$ is the physical matching scale of the construction depends on
information not yet in the model (the quadratic divergence is regulator
dependent and is not defined in $\overline{\rm MS}$). An alternative
EFT description --- a dim-6 $(H^\dagger H)^3 / \Lambda_d^2$ operator
with $c_6 = 1$ at $\Lambda_d \simeq 1.15$~TeV --- reproduces $\Delta_d$
at the same order, with the Wilson coefficient remaining a free
parameter.

\paragraph{Mechanism for the triplet shift: 15$\times$ shortfall at one loop.}
Producing the triplet half of the bidirectional shift is a distinct
theoretical question. The natural candidate -- one-loop SU(2)$_L$
gauge-boson contributions to the triplet's Coleman--Weinberg potential
-- has the correct \emph{sign} (gauge loops on a tachyonic scalar
make $m^2$ less negative, lowering $|\lambda_{\mt,-}|m_0^2$ as
required). Computed at the natural cutoff $\Lambda = m_0$
(\texttt{lagrangian/Q6\_triplet\_gauge\_CW.py}):
\begin{equation}
  \Delta M_{\mt}^2(\text{gauge}) \;=\;
    \frac{3 g^2}{16 \pi^2} \cdot \Cas(\mt) \cdot \frac{m_0^2}{4}
    \;\simeq\; (6.78~\text{GeV})^2.
\end{equation}
This is $\sim\!15\times$ smaller than the required $\Delta_t = (26.77~\text{GeV})^2$:
the gauge-loop mechanism alone gives the right direction but not the
right magnitude. The triplet has no Yukawa coupling to SM fermions
(Y=0 triplet), so the analogue of the doublet's top-CW contribution
($\sim\!(26.4~\text{GeV})^2$) is absent.

The bidirectional spurion picture's triplet half therefore lacks a
natural one-loop derivation, and we honestly distinguish two readings:

\begin{itemize}
\item (\emph{Empirical:}) the cross-irrep negative-branch trace
  identity Eq.~\eqref{eq:trace_id} holds to $(5~\text{GeV})^2$ as a
  numerical coincidence within the model, with the doublet shift
  derived (top-CW) and the triplet shift unexplained at one loop.
\item (\emph{Mechanistic, hypothetical:}) some additional UV
  contribution (heavy composite-quark loops via specific PC
  Yukawas; two-loop effects; or a portal-induced back-reaction)
  produces the missing $\sim\!(26~\text{GeV})^2$ on the triplet.
  None of the obvious candidates is quantitatively established.
\end{itemize}

In the empirical reading, $v/\sqrt{2}$ is matched to $0.04\%$ via the
trace identity but is not strictly a model output; the genuine
parameter-free outputs are $M_Z$ (from the gauge sector) and
$m_{h,\rm bare}$ (from the doublet seed eigenvalue, Q2 below).

\paragraph{Special case: unidirectional spurion.} The earlier draft of
this paper used only the doublet half of Eq.~\eqref{eq:bidir} and
treated $|\lambda_{\mt,-}|\,m_0^2$ as a separate unidentified heavy
state at $176$~GeV. That picture is recovered as the special case
where the triplet sector is decoupled from the spurion mechanism;
the bidirectional ansatz subsumes it and adds a prediction (the
Fermi scale).

\section{MadGraph consistency check}

A hand-written UFO file (built on the bundled FR-SM via
\texttt{model/ufo/setup\_ufo.sh}) substitutes the Casimir-locked
Higgs-mass value $m_h = 125.30$~GeV for the SM default of $125.0$~GeV.
Imported into MadGraph5\_aMC@NLO~3.7.1, the resulting model yields
the following LO cross-sections at LHC $13$~TeV:
\begin{center}
\begin{tabular}{l c c}
\toprule
Process & $\sigma_{\rm LO}$ (this model, pb) & $\sigma_{\rm LO}$ (SM benchmark, pb)\\
\midrule
$p p \to W^+ W^-$ & 64.46 & $\sim 65{-}70$ \\
$p p \to Z Z$     & 9.23  & $\sim 9{-}11$  \\
$p p \to H Z$     & 0.578 & $\sim 0.55{-}0.62$ \\
\bottomrule
\end{tabular}
\end{center}
We emphasise that these are \emph{not} non-trivial validations of the
Casimir-locking hypothesis. The tree-level diboson amplitudes are
unchanged from the SM (no new vertices were added), so SM-like rates are
the \emph{expected and required} outcome --- if they had deviated, the
UFO would have a bug. The MG5 check is a necessary technical consistency
test that the UFO compiles and reproduces SM IR phenomenology at the
$m_h = 125.30$~GeV input point. Non-trivial collider tests require
adding the BSM scalars as new fields (deferred to future work), so that
their production and decay can be computed.

\section{Electroweak precision-test estimate}

We \emph{estimate} the Peskin--Takeuchi~\cite{PeskinTakeuchi1992}
contributions $\Delta S$, $\Delta T$, $\Delta U$ from the $Y = 0$ triplet
using a leading-log approximation for $\Delta S$ and a closed-form
scalar loop function for $\Delta T$, following the literature on
$Y=0$ triplet scalars~\cite{ForshawRossSankey1991, Khan2018Triplet, Katayose2021TripletDM}. In the
custodial limit (degenerate triplet components) all three vanish.
For a small custodial-breaking splitting $\Delta M = 5$~GeV between
neutral and charged components, the leading-log estimate gives
\(\Delta S \simeq -2.8\times 10^{-3}\),
\(\Delta T \simeq +2.3\times 10^{-4}\),
\(\Delta U \simeq +8\times 10^{-6}\),
all far inside the PDG~2024 ranges
($S = -0.04 \pm 0.10$, $T = 0.01 \pm 0.12$, $U = -0.01 \pm 0.09$).
A full one-loop Passarino--Veltman evaluation would be required for a
publication-grade precision claim; we have not performed that here. The
present estimate is a consistency check that the model is not grossly
incompatible with EWPT, not a precision phenomenology result. Future
work using HEPfit or smelli is required for the proper $\chi^2$ fit.

\section{The j=3/2 channel: a structural side-prediction}
\label{sec:j32}

The Casimir-locked seed is naturally extensible to internal SU(2)
representation labels beyond $\md$ and $\mt$. For each irrep
$R = 2j+1$, the same seed structure
$\mathcal M_R^2 / m_0^2 = \big[\begin{smallmatrix} 0 & \sqrt{\Cas(R)} \\ \sqrt{\Cas(R)} & -\Cas(R)\end{smallmatrix}\big]$
gives two physical scalar masses
\begin{equation}
  M_{R,\pm}^2 = m_0^2\,\frac{-\Cas(R) \pm \sqrt{\Cas(R)^2 + 4\Cas(R)}}{2}.
\end{equation}
For the next irrep, $R = 4$ (internal label $j = 3/2$), $\Cas(4) = 15/4$
and the closed form is
\begin{equation}
  \lambda_{R=4,+} = \frac{-15 + \sqrt{465}}{8},\qquad
  M(\Phi^4_+) \;=\; m_0\,\sqrt{\lambda_{R=4,+}}
  \;=\; 96.54~\text{GeV}.
\label{eq:96GeV}
\end{equation}
With $m_0$ fixed by the gauge anchor (Sec.~6) there is no remaining free
parameter, so Eq.~\eqref{eq:96GeV} is a parameter-free output of the
Casimir-locked structure. Table~\ref{tab:tower} summarises the spectrum
tower.

\begin{table}[ht]
\centering
\begin{tabular}{c c c c c l}
\toprule
$R$ & $j$ & $\Cas(R)$ & $M_{R,+}$ (GeV) & $M_{R,-}$ (GeV) & content \\
\midrule
2 & 1/2 & 3/4   & 80.37  & 122.39 & doublet seed; $M_{R,+} = M_W$ \\
3 & 1   & 2     & 91.19  & 176.16 & triplet seed; $M_{R,+} = M_Z$ \\
4 & 3/2 & 15/4  & 96.54  & 227.85 & v2 prediction (no observable identified) \\
5 & 2   & 6     & 99.58  & 279.41 & not in scope \\
\bottomrule
\end{tabular}
\caption{Casimir-locked spectrum tower. Each internal SU(2) irrep $R$
contributes one positive- and one negative-eigenvalue scalar.}
\label{tab:tower}
\end{table}

\paragraph{Comment on the CMS+LEP $\gamma\gamma$ hint.} The
$96.54$~GeV value sits in the mass range of the CMS+LEP $\gamma\gamma$
excess near $95.4$~GeV. We do \emph{not} claim a quantitative match:
the deviation is $1.2\,\%$, roughly $30\times$ looser than the
$\{M_W, M_Z, m_h, v/\sqrt{2}\}$ predictions of Sections~6 and~7, which
hold to $0.04\,\%$ or better. We list it as a structural side-prediction
that lands in a window of current experimental interest, not as a
confirmation of the model.

\paragraph{Hans's spin-3/2.} A natural objection (chat sec.~10) was that
the model labels look like physical Lorentz spins; in particular, an
internal label $j = 3/2$ would suggest a Rarita--Schwinger fermion in the
spectrum. The chat dismissed that reading: the labels are
\emph{representation indices}, not Lorentz spins. The $R = 4$ result
above \emph{rescues} the spin-3/2 intuition without giving up the
dismissal: the predicted state is a \emph{scalar}, but its mass is fixed
by the $j = 3/2$ Casimir of an internal channel. The spin-3/2 ``label''
lives on as a representation tag, not as a propagating spin.

\paragraph{Phenomenology of the $R=4$ multiplet.} The 4-dimensional
SU(2) irrep contains $T_3 \in \{-3/2, -1/2, +1/2, +3/2\}$. With
hypercharge $Y = 1/2$ the electric charges are $Q \in \{-1, 0, +1, +2\}$,
including a doubly-charged scalar $\Phi^{++}$ at a mass close to
$M(\Phi^4_+) = 96.5$~GeV. Existing LHC searches for doubly-charged
scalars (ATLAS, CMS) constrain the same multiplet and would be the
sharpest direct test if the $R = 4$ channel is realised.

\section{Phenomenology of the 91~GeV neutral scalar}
\label{sec:91GeV}

The model's most distinctive prediction is the neutral scalar
$\Phi^t_+$ at exactly $M_Z = 91.188$~GeV (Section~6, Distinctive triplet
prediction). This section addresses Open Question Q4: \emph{is this state
already excluded by LEP and LHC?} A detailed numerical recast is in
\texttt{phenomenology/recast\_91GeV.py}.

\paragraph{Structural couplings.} As the positive mass eigenstate of the
$Y=0$ SU(2)$_L$ triplet sector, $\Phi^t_+$ has:
\begin{itemize}
\item \emph{No tree-level Yukawa couplings to SM fermions:} the operator
  $\bar L_L\, \Phi^t\, e_R$ is gauge-allowed only for an SU(2) doublet
  scalar with $Y = 1/2$, not for a $Y=0$ triplet.
\item \emph{No direct $Z\Phi^t_+\Phi^t_+$ vertex:} the triplet covariant
  derivative gives only $g\, T^a (D Z)^b (D Z)^c \epsilon^{abc}$, which
  does not contain a neutral-neutral $\Phi^t$ pair.
\item \emph{Higgs portal $\lambda_{dt}$ generates production} through
  the doublet--triplet mixing $\sin^2\theta_{\rm mix} \sim \lambda_{dt}\,v^2 / (M_h^2 - M_{\Phi^t_+}^2)$
  $\simeq 12.6\,\lambda_{dt}$.
\end{itemize}
$\Phi^t_+$ is therefore both gauge-phobic and fermion-phobic in the bare
limit; observability is controlled entirely by $\lambda_{dt}$.

\paragraph{LEP HZ$^\prime$ constraint.} The combined
ALEPH+DELPHI+L3+OPAL Higgs-strahlung search (LHWG-Note-2005-01) yields
$\sigma(e^+e^- \to Z h^\prime)/\sigma_{\rm SM}(m=91~{\rm GeV}) < 0.15$
in the $91$~GeV window, hence
$\sin^2\theta_{\rm mix} < 0.15$, i.e.\
\begin{equation}
  \lambda_{dt} \;<\; 0.012
  \qquad (\text{LEP HZ}^\prime, \text{combined}).
\end{equation}
This is the strongest current constraint and is automatically satisfied
by the $\rho$-parameter requirement from Section~5 that the triplet
sector remain custodially-protected (any $\lambda_{dt}$ exceeding the
LEP bound would also produce a triplet VEV exceeding the
$\rho$-parameter limit).

\paragraph{LHC diphoton.} The ATLAS search for a low-mass scalar in the
$65$--$110$~GeV window~\cite{Khan2018Triplet, Katayose2021TripletDM}
sets $\sigma(pp \to h^\prime) \times \mathrm{BR}(h^\prime \to
\gamma\gamma) < 100$~fb at $13$~TeV. With SM-Higgs-like gluon-fusion
production (${\sim}80$~pb at $91$~GeV) and SM-Higgs-like diphoton
branching ratio ($\sim 1.4\times 10^{-3}$ at $91$~GeV), the SM-strength
signal is $\sim 112$~fb, just at the edge of the experimental bound.
This translates to $\sin^2\theta_{\rm mix} < 0.89$, weaker than the LEP
bound by an order of magnitude.

\paragraph{Verdict.} The 91~GeV neutral scalar is \emph{not} excluded by
current data. The LEP HZ$^\prime$ bound $\lambda_{dt} < 0.012$ is
already required by the $\rho$-parameter constraint, so the model
needs no additional tuning to evade existing limits. A future $e^+e^-$
Higgs factory at $\sqrt s = 240$~GeV (FCC-ee, CEPC) would probe
$\sin^2\theta_{\rm mix}$ to $\sim 10^{-3}$ and $\lambda_{dt}$ to
$\sim 10^{-4}$ -- a definitive direct test of this prediction. Open
Question Q4 of Section~\ref{sec:status} is therefore answered
affirmatively: the model survives.

\section{Status of the derivation and open questions}
\label{sec:status}

We summarise here, transparently, what this draft does and does not show.

\paragraph{What is in the paper.}
\begin{enumerate}
\item A scalar field content (doublet + triplet, each in an
  elementary/composite basis) and a quadratic seed potential with
  Casimir-locked entries [Eq.~\eqref{eq:seed}].
\item Closed-form diagonalisation of the seed matrices and the
  closed-form four-eigenvalue spectrum [Eq.~\eqref{eq:eigs}].
\item Two \emph{postulated} identifications mapping the seed eigenvalues
  onto physical masses [Eqs.~\eqref{eq:gauge_id} and~\eqref{eq:higgs_bare}].
\item The observation that, conditioned on those postulates, the
  numerical values of $M_W$ and $m_h$ are reproduced to $10^{-3}$ given
  $M_Z$ and $G_F$ as external inputs.
\item A parameter-free side-prediction from extending to $j=3/2$
  (Eq.~\eqref{eq:96GeV}): a scalar at $96.54$~GeV. This sits in the
  CMS+LEP $\gamma\gamma$ window near $95.4$~GeV but at $1.2\,\%$
  deviation; we do not claim the identification.
\item A consistency check that the modified UFO ($m_h = 125.30$~GeV)
  reproduces SM-like LO diboson cross-sections in MadGraph
  (necessary but not non-trivial).
\item A leading-log estimate of $\Delta S$, $\Delta T$, $\Delta U$
  showing the model is not grossly incompatible with EWPT.
\end{enumerate}

\paragraph{What is open work.}
\begin{description}
\item[Q1.] (\emph{Partially answered in Sec.~\ref{sec:PCsketch}.}) The
  explicit chain UV $\to$ Casimir-locked seed $\to$ EWSB $\to$ gauge
  mass matrix CLOSES into a consistent SO(5)/SO(4) composite-Higgs
  framework (\texttt{lagrangian/Q1\_full\_chain.py}). The
  identifications $M_W^2 = m_0^2 \lambda_{\md,+}$ and
  $M_Z^2 = m_0^2 \lambda_{\mt,+}$ are equivalent (via standard EWSB on
  the doublet's tachyonic eigenstate) to algebraic relations between
  $g$, $g'$, $v_H$, and $m_0$ that are satisfied identically with
  $m_0$ anchored by $M_W$ and $v_H$ from the SO(5)/SO(4) Coleman--
  Weinberg minimum. The non-trivial PHYSICS prediction extracted from
  the closure is $\sin^2\theta_W = 1 - \lambda_{\md,+}/\lambda_{\mt,+}$,
  parameter-free. \emph{Still open}: a quantitative UV derivation
  showing the CW minimum lands on $v_H = 246.22$~GeV with $m_0 =
  106.57$~GeV (requires fixing $N_{\rm HC}$, fermion content, Yukawa
  structure -- multi-week model-building).
\item[Q2.] (\emph{Largely answered in
  \texttt{lagrangian/Q2\_potential\_minimisation.py}.}) Under the
  canonical real-scalar normalisation $V = \tfrac{1}{2} m^2 |\Phi|^2 +
  (\lambda_h/4) |\Phi|^4$, the bare Higgs mass-squared is fixed by the
  TACHYONIC EIGENVALUE OF THE SEED ALONE: $m_{h,\rm bare}^2 = m_0^2
  |\lambda_-(2)|$, with NO contribution from the quartic couplings at
  tree level. The quartics determine the Higgs VEV $v$ via an effective
  combination $\lambda_h^{\rm eff} = \lambda_{ee} p_e^2 + \lambda_{cc}
  p_c^2 + \lambda_{ec}\, p_e p_c + \lambda_{ec}'\, (c_e^- c_c^-)^2$,
  weighted by the mixing-angle projections. For $v_{\rm SM}=246.22$~GeV
  this requires $\lambda_h^{\rm eff} = 0.247$ -- one equation for four
  bare quartics, leaving a 3-parameter family of viable Lagrangians.
  Whether the partial-compositeness UV uniquely fixes $\lambda_h^{\rm
  eff}$ remains the natural follow-up; the standard composite-Higgs
  expectation $\lambda_h \sim N_c y_t^4 / 16\pi^2 \sim 0.06$ is $\sim
  4\times$ smaller, a manifestation of the well-known ``small
  $\lambda_h$ problem'' of composite Higgs.
\item[Q3.] Which definition of $\sin^2\theta_W$ is being matched? The
  algebraic value $0.22310$ from Eq.~\eqref{eq:sin2thW} should be
  compared against the on-shell scheme value
  $\sin^2\theta_W^{\rm OS} = 1 - M_W^2/M_Z^2 = 0.22290(28)$. The
  agreement is at the $1\sigma$ level, not at $10^{-10}$.
\item[Q4.] (\emph{Answered in Sec.~\ref{sec:91GeV}.}) The 91.2~GeV
  neutral scalar $\Phi^t_+$ has no tree-level Yukawa couplings and no
  direct $Z\Phi^t_+\Phi^t_+$ vertex; it is fermion- and gauge-phobic.
  The LEP HZ$^\prime$ combined search constrains
  $\lambda_{dt} < 0.012$, automatically satisfied by the
  $\rho$-parameter requirement. The state is \emph{not} excluded by
  current data; FCC-ee/CEPC at $\sqrt{s} = 240$~GeV would probe it
  definitively.
\item[Q5.] Are Eqs.~\eqref{eq:gauge_id} and~\eqref{eq:higgs_bare} a
  consequence of a UV-complete partial-compositeness or top-seesaw
  construction (\textit{cf.}~\cite{Caracciolo2012, Chivukula1998}), or
  do they require a new dynamical principle?
\item[Q6.] (\emph{Partially answered in
  \texttt{lagrangian/Q6\_triplet\_gauge\_CW.py}; sign correct,
  magnitude $15\times$ short.}) The triplet's one-loop gauge
  contribution at $\Lambda = m_0$ has the right SIGN (gauge loops on
  a tachyonic scalar reduce $|m^2|$, lowering $|\lambda_-(\mt)|m_0^2$
  toward $(v/\sqrt{2})^2$ as required), but the magnitude $\Delta
  M_\mt^2(\text{gauge}) = (3 g^2/16\pi^2) \Cas(\mt) m_0^2 / 4 = (6.78~
  \text{GeV})^2$ is $15\times$ smaller than the required $\Delta_t =
  (26.77~\text{GeV})^2$. The triplet has no Yukawa to SM fermions
  ($Y=0$), so the analog of the doublet's top-CW contribution is
  ABSENT. The bidirectional picture's triplet half therefore lacks a
  natural one-loop derivation: in the empirical reading
  (Sec.~\ref{sec:spurion}) the trace identity Eq.~\eqref{eq:trace_id}
  is a $(5~\text{GeV})^2$ numerical coincidence with the doublet
  shift derived (top-CW) and the triplet shift unexplained at one
  loop. The genuine parameter-free outputs of the model are $M_Z$
  and $m_{h,\rm bare}$; $v/\sqrt{2}$ is matched but not strictly
  derived.
\end{description}

In summary, the current draft is best read as a \emph{numerically
suggestive ansatz} for the electroweak boson spectrum, with $M_W$,
$m_h$, $v/\sqrt{2}$ as the testable numerical consequences (all
matched to $0.04\,\%$ or better) and a $96.5$~GeV scalar as a
parameter-free side-prediction at lower precision. Promoting it to a
derived model requires answering Q1--Q6 above.

\section{Sketch of a partial-compositeness origin for the $a = b$ lock}
\label{sec:PCsketch}

Phase A established that the $a = b$ Casimir lock is the operational
content of the model. We sketch here a UV completion in which $a = b$
emerges naturally; full derivation is multi-week work, deferred.

\paragraph{Setup} (Caracciolo--Parolini--Serone-style~\cite{Caracciolo2012}):
\begin{itemize}
\item Hypercolor gauge group SU($N$)$_{\rm HC}$, strongly coupled at
  scale $\Lambda_{\rm HC}\!\sim\!$ TeV.
\item SU(2)$_L$ gauged as a subgroup of a global symmetry of the HC sector.
\item Composite scalar $\chi_R$: bound state of HC fermions, in the
  irrep $R$ of SU(2)$_L$.
\item Elementary scalar $\psi_R$: SM-like partner with protected
  diagonal zero (the $(1,1)$ entry).
\item Mixing operator generated by an adjoint-charged spurion
  $\epsilon^A$:
  \begin{equation}
    \mathcal L_{\rm mix} \;=\; y\,\epsilon^A\,\psi_R^a\,(T_R^A)_{ab}\,\chi_R^b + \text{h.c.}
  \label{eq:Lmix}
  \end{equation}
\end{itemize}

\paragraph{Result.} The off-diagonal mixing in the $(\psi, \chi)$ basis,
upon projecting Eq.~\eqref{eq:Lmix} onto the $\Cas(R)$ component using
$\langle T^A T^A\rangle_R = \Cas(R)\,\mathbf{1}_R$, evaluates to
\begin{equation}
  M_{\rm mix} \;=\; y\,|\epsilon|\,\sqrt{\Cas(R)}.
\end{equation}
Comparing with the Casimir-locked seed $a\,\sqrt{\Cas(R)}\,m_0^2$ gives
$a\,m_0^2 = y\,|\epsilon|$. The composite mass-squared in the
$(\chi)$ entry is $-m_\chi^2$, comparing with $-b\,\Cas(R)\,m_0^2$ gives
$b\,m_0^2 = m_\chi^2 / \Cas(R)$. Therefore the Casimir lock $a = b$
holds in this UV completion if and only if
\begin{equation}
  m_\chi^2 \;=\; y\,|\epsilon|\,\Cas(R),
\label{eq:ab_condition}
\end{equation}
i.e., the composite mass-squared is proportional to the same coupling
$y\,|\epsilon|$ that generates the elementary--composite mixing, with
prefactor $\Cas(R)$. This is a strong constraint on the HC dynamics, but
physically motivated: it holds if both $m_\chi^2$ and the mixing
descend from the \emph{same} higher-dimensional operator above
$\Lambda_{\rm HC}$, a generic feature of partial-compositeness
constructions.

\paragraph{Conceptual amusement.} It is worth noting that the chat
draft (Section~10) dismissed the reading $j = $ Lorentz spin --- the
labels are internal SU(2)$_L$ representation indices, not propagating
spin --- and the partial-compositeness origin sketched above
\emph{vindicates} that dismissal in a structurally interesting way: the
$j$ labels are the SU(2)$_L$ charges of \emph{composite scalar bound
states} (the $\chi_R$). They are structurally analogous to the isospin
labels of QCD mesons (a pion is a $J=0$ scalar with $I=1$, a kaon is
$J=0$ with $I=1/2$, etc.). The Casimir-locked spectrum is the analogue
of an isospin tower of hypercolor mesons. Same group theory; different
physical referent.

\subsection*{The full chain UV $\to$ $M_W^2, M_Z^2$ (closure check)}

We can chain the partial-compositeness sketch above with the
SO(5)/SO(4) Coleman--Weinberg potential and standard EWSB to check
\emph{internal consistency} of the Casimir-locked identification with a
fully-derived Lagrangian-level Higgs sector. The full chain is
implemented in \texttt{lagrangian/Q1\_full\_chain.py} and proceeds
in five stages:

\begin{enumerate}
\item \emph{UV Lagrangian}: SU($N$)$_{\rm HC}$ + SO(5)/SO(4) coset +
   partial-compositeness mixings $y\,\epsilon^A \psi_R T^A_R \chi_R$.
\item \emph{Casimir lock}: PC mixings produce the off-diagonal
   $y\,|\epsilon|\,\sqrt{\Cas(R)}$; the lock $a = b$ requires
   $m_\chi^2 = y\,|\epsilon|\,\Cas(R)$ (Eq.~\eqref{eq:ab_condition}),
   yielding the canonical seed Eq.~\eqref{eq:seed} with
   $m_0^2 = y\,|\epsilon|$.
\item \emph{Doublet sector EWSB}: in canonical scalar normalisation
   $V = \tfrac{1}{2} m^2 |\Phi|^2 + (\lambda_h/4)|\Phi|^4$, the seed
   eigenvalue $m_0^2 \lambda_-(2)$ \emph{is} the bare physical
   Higgs mass-squared, no factor-of-2 shift. EWSB on the doublet's
   tachyonic mode $\Phi_-^d$ produces $v_H = 246.22$~GeV after
   the spurion-shifted potential is minimised.
\item \emph{Coleman--Weinberg minimisation}: the SO(5)/SO(4) potential
   $V(h) = a \sin^2(h/f) - b \sin^2(h/f)\cos^2(h/f)$ has minimum
   $y_* = \sin^2(h/f) = (b-a)/(2b)$, $v_H = f\,\sqrt{y_*}$, with
   $a, b$ from one-loop gauge and top contributions.
\item \emph{Gauge mass matrix}: EWSB on $\langle h_0\rangle = v_H/\sqrt{2}$
   gives the standard
   \begin{equation}
     M_W^2 = \frac{g^2 v_H^2}{4},\qquad
     M_Z^2 = \frac{(g^2 + g'^2)\,v_H^2}{4}.
   \label{eq:gauge_from_vH}
   \end{equation}
\end{enumerate}

\paragraph{Closure.} The Casimir-locked identification of
Section~\ref{sec:postulates}, $M_W^2 = m_0^2\,\lambda_+(2)$ and
$M_Z^2 = m_0^2\,\lambda_+(3)$, is then equivalent (via
Eq.~\eqref{eq:gauge_from_vH}) to two algebraic relations:
\begin{equation}
  g^2 \;=\; \frac{4\,m_0^2\,\lambda_+(2)}{v_H^2},\qquad
  g'^2 \;=\; \frac{4\,m_0^2\,(\lambda_+(3) - \lambda_+(2))}{v_H^2}.
\end{equation}
With $m_0$ anchored by $M_W$ (Section~\ref{sec:postulates}) and
$v_H$ from the SO(5)/SO(4) CW minimum, both relations are satisfied
identically by construction. The non-trivial physics output extracted
from the closure is the parameter-free Casimir-locked Weinberg angle:
\begin{equation}
  \sin^2\theta_W \;=\; 1 \;-\; \lambda_+(2)/\lambda_+(3) \;=\; 0.22310,
\end{equation}
matching PDG$_{\rm OS}$ to $0.09\%$.

\paragraph{What the chain shows, and what it does not.}
The chain demonstrates that the Casimir-locked identification
\emph{closes into} a fully-derived SO(5)/SO(4) composite-Higgs
framework. It is \emph{not} a first-principles derivation of the
specific value $m_0 = 106.57$~GeV from the UV: that requires fixing
$N_{\rm HC}$, the strong-sector fermion content, and the Yukawa
structure such that the CW minimum lands on $v_H = 246.22$~GeV with
the correct relation to $m_0$. Showing the SO(5)/SO(4) CW minimum
quantitatively reproduces the seed eigenvalues is multi-week
UV model-building, deferred to follow-up work. The framework
hosts the Casimir lock without inconsistency; the construction
is now of the same theoretical status as the classic composite-Higgs
constructions, with the additional Casimir-locking constraint
that picks out the doublet--triplet pair (Sec.~\ref{sec:repsel}) as
the simplest realisation.

\section{Discussion and outlook}

The Casimir-locked ansatz is a deformation of known BSM physics
(bosonic seesaw + partial compositeness + weak triplet) by a
representation-theoretic constraint on its mass-squared matrix entries.
Conditioned on the postulates of Sections~6--7, the framework reproduces
$\{M_W, M_Z, m_h, v/\sqrt{2}\}$ at the $10^{-3}$--$10^{-4}$ level using
$M_W$ as the single anchor and one spurion magnitude $\Delta$, and
yields a $96.5$~GeV scalar from a single additional irrep ($j = 3/2$)
as a parameter-free side-prediction.

The most pressing extensions, deferred to future work, are: (i) embedding
the Casimir-locked seed into a UV-complete partial-compositeness
framework~\cite{Caracciolo2012} or top-seesaw
construction~\cite{Chivukula1998} that generates the seed dynamically;
(ii) extending the analysis to the fermion sector via the same
elementary/composite basis; and (iii) collider phenomenology of the
$M = 91.2$~GeV neutral scalar predicted by the triplet positive
eigenvalue.

\paragraph{Remark on a left-right symmetric extension.} A natural
question raised in the present work is whether U(1)$_Y$ itself can be
viewed as descending from a Casimir-locked seed structure on SU(2)$_R$.
In the standard left-right symmetric framework
$SU(2)_L \times SU(2)_R \times U(1)_{B-L}$ with hypercharge
$Y = T^3_R + (B-L)/2$, an SU(2)$_R$ Casimir-locked seed with its own
scale $m_0^R$ would generate a triplet $\Delta_R$ whose tachyonic
eigenvalue breaks $SU(2)_R \times U(1)_{B-L} \to U(1)_Y$. The
parameter-free testable prediction of this realisation is
\begin{equation}
  \frac{M(W_R)}{M(Z_R)} \;=\; \sqrt{\frac{\lambda_+(\md)}{\lambda_+(\mt)}} \;=\; 0.881,
\end{equation}
i.e., the heavy left-right partners would share the SAME mass ratio
as $M_W/M_Z$ in the low-energy sector, independent of $m_0^R$. Strict
L-R parity ($m_0^R = m_0^L$) is excluded by current LHC bounds
$M(W_R) > 4$--$6$~TeV, so the two seed scales must be independent;
this trades one input ($g'$) for two ($m_0^R$, $g_{B-L}$) and increases
the parameter count overall. We therefore present this as a future
direction rather than an in-line proposal; the standalone analysis is
in \texttt{lagrangian/LRSM\_extension\_mock.py}.

\section*{Computational pipeline}

All numerical results in this paper are reproducible end-to-end via a
public pipeline:
\begin{itemize}
\item \texttt{algebra/seeds.py} --- Casimir-locked seed matrices and
   eigenvalue derivation in SymPy.
\item \texttt{algebra/ewsb\_check.py} --- end-to-end mass derivation;
   exits with code 0 iff every prediction matches PDG to better than
   $0.5\%$.
\item \texttt{lagrangian/L\_gauge.py}, \texttt{lagrangian/L\_scalar.py}
   --- symbolic Lagrangian.
\item \texttt{feynrules/casimir\_seesaw.fr} --- FeynRules 2.x model file.
\item \texttt{madgraph/checks.mg5} --- MG5\_aMC@NLO process cards for
   tree-level diboson cross-section sanity checks.
\item \texttt{ewpt/STU.py} --- oblique parameters with PDG confrontation.
\item \texttt{run\_all.sh} --- one-line invocation that runs the entire
   pipeline.
\end{itemize}

\section*{Acknowledgements}

We thank the participants of the BSM-alternatives-to-SU(2)$\times$U(1)
discussion for the structural framing of the Casimir lock and the
historical context of Higgs's ``incomplete multiplets'' remark.

\paragraph{Use of generative AI.} In accordance with Physical Review D
editorial policy on the use of artificial intelligence in scientific
manuscripts, we disclose that portions of this work --- including
draft preparation, symbolic algebra (SymPy), numerical scans
(\texttt{algebra/}, \texttt{phenomenology/}), and the assembly of
the literature comparison table --- were carried out with assistance
from the Claude AI assistant (Anthropic). All physical interpretations,
derivations, identifications, and numerical results were independently
verified by the authors. The authors take full responsibility for the
content of this manuscript. The complete computational pipeline
(\texttt{run\_all.sh}) reproduces every numerical statement in this
paper end-to-end.

\bibliographystyle{plain}
\bibliography{../../refs/refs}

\end{document}
